\scanpage{4}

$E$ faisceau loc. libre de rang $r+1$ sur $S$

$P=P(E)$ fibré projectif associé, on a une suite exacte canonique
\diagram{\begin{tikzcd}
0\arrow[r]&\Omega^1_{P/S}\arrow[r]&g^*(\underline E)(-1)\arrow[r,dash]&\underline O_P\arrow[r,dash]&0
\end{tikzcd}}

[qui donne un isom can
\[
(1)\qquad\Omega^r_{P/S}\xrightarrow{\sim}g^*\Bigl(\bigwedge_{\underline O_S}^{r+1}E\Bigr)(-r-1)
\qquad]
\]

Si donc $X$ est un sous-schéma de $P$ \unsure{lisse} sur $S$, de dim relative $r-p$
\inserted{$(n=$}
associé au faisceau \deleted{loc. libre} \inserted{d'idéaux} $J$, on a la suite exacte
\diagram{\begin{tikzcd}
0\arrow[r,dash]&J/J^2\arrow[r]&\Omega^1_{P/S}\otimes\underline O_X\arrow[r]&\Omega^1_{X/S}\arrow[r]&0
\end{tikzcd}}
d'où l'isomorphisme
\[
(2)\qquad\Omega^r_{P/S}\otimes\underline O_X
\xleftarrow{\sim}\Omega^{\unsure{n}}_{X/S}\otimes\bigwedge_{\underline O_X}^p(J/J^2)
\]
d'où en composant (1) et (2), désignant par $f:X\to S$
le morphisme structural
\[
\deleted{$\unsure{\Omega^r}$}
f^*\Bigl(\bigwedge_{\underline O_S}^{r+1}\underline E\Bigr)(-r-1)
\xleftarrow{\sim}\Omega^{\unsure{n}}_{X/S}\otimes\bigwedge_{\underline O_X}^p(J/J^2)
\]

Supposons que $X$ soit une intersection complète, définie
par un sous-module \inserted{gradué $M$} loc. facteur direct de
$\operatorname{Symm}^*(\underline E)$,
\deleted{[qui est \unsure{un}} loc. libre de rang $p$
[$\unsure{\text{il n'est}}$ \deleted{défini}
déterminé que modulo le carré $\underline I^2$ de l'Idéal qu'il
définit ; mais sa structure est canonique]. On a
un hom., si $I_d$ est la composante de degré $d$ de $I$
\[
g^*(I_d)(-d)\longrightarrow J
\]
d'où des hom \uline{surjectifs}
\[
\coprod g^*(M_d)(-d)\longrightarrow J
\]
induisant un \uline{isom}
\[
\coprod f^*(M_d)(-d)\xrightarrow{\sim}J/J^2\simeq\operatorname{Conormal}_{P/X}
\]

\mspage{3}
